\section{Security Considerations}
\label{sec:security}
%% ============================================================

\subsection{Encryption Layers}

The system provides defense in depth through three independent
encryption layers:

\begin{table}[H]
\centering
\caption{Encryption layers.}
\label{tab:encryption}
\begin{tabular}{@{}lllp{5cm}@{}}
\toprule
\textbf{Layer} & \textbf{Scope} & \textbf{Algorithm} & \textbf{Purpose} \\
\midrule
Application & On-disk DB & SQLCipher AES-256-CBC & Protects data at rest
  on pod local disk. Key derived from KMS. \\
Replication & S3 objects & age X25519 + ChaCha20-Poly1305 & Protects
  data at rest in object storage. Per-org keys. \\
Transport & Network & TLS 1.3 & Protects data in transit to S3. \\
\bottomrule
\end{tabular}
\end{table}

An attacker who compromises S3 sees only age-encrypted blobs. An
attacker who compromises the pod filesystem sees only
SQLCipher-encrypted databases. An attacker who compromises the network
sees only TLS-encrypted traffic. All three layers must be breached
simultaneously to access plaintext.

\subsection{Key Hierarchy}\label{sec:key-hierarchy}

\begin{figure}[H]
\centering
\begin{tikzpicture}[
  key/.style={draw, rounded corners=2pt, minimum width=3cm,
    minimum height=0.7cm, font=\small},
  arr/.style={-{Stealth[length=4pt]}, thick},
]
  \node[key, fill=red!10] (master) at (0, 3)
    {Master Key (KMS)};
  \node[key, fill=blue!10] (orgA) at (-3, 1.5)
    {Org A Key (HKDF)};
  \node[key, fill=blue!10] (orgB) at (0, 1.5)
    {Org B Key (HKDF)};
  \node[key, fill=blue!10] (orgN) at (3, 1.5)
    {Org N Key (HKDF)};
  \node[key, fill=green!10] (ageA) at (-3, 0)
    {age Identity A};
  \node[key, fill=green!10] (ageB) at (0, 0)
    {age Identity B};
  \node[key, fill=green!10] (ageN) at (3, 0)
    {age Identity N};
  \node[key, fill=yellow!10] (recovery) at (5.5, 3)
    {Recovery Key (KMS)};

  \draw[arr] (master) -- (orgA);
  \draw[arr] (master) -- (orgB);
  \draw[arr] (master) -- (orgN);
  \draw[arr] (orgA) -- (ageA);
  \draw[arr] (orgB) -- (ageB);
  \draw[arr] (orgN) -- (ageN);
\end{tikzpicture}
\caption{Key derivation hierarchy. The master key in KMS derives
  per-org keys via HKDF. Each org key generates an age identity.
  The recovery key is independent and added as a second recipient.}
\label{fig:keyhier}
\end{figure}

\begin{invariant}[No Plaintext Egress]\label{inv:no-plaintext}
No plaintext database content ever exists outside the pod boundary.
WAL frames are encrypted before upload. Snapshots are encrypted before
upload. The age identity (private key) exists only in the pod's
memory, loaded from a K8s Secret mounted as a tmpfs volume.
\end{invariant}

\subsection{Threat Model}

\begin{table}[H]
\centering
\caption{Threat model and mitigations.}
\label{tab:threats}
\begin{tabular}{@{}lp{4cm}p{5cm}@{}}
\toprule
\textbf{Threat} & \textbf{Attack Vector} & \textbf{Mitigation} \\
\midrule
S3 compromise & Attacker gains read access to S3 bucket &
  All objects age-encrypted; no plaintext. \\
Pod compromise & Attacker gains shell in pod &
  SQLCipher encrypts on-disk DB; age identity in tmpfs. \\
KMS compromise & Attacker reads master key &
  All org keys derivable; requires KMS + S3 compromise
  simultaneously for data access. \\
Rogue operator & Operator uses recovery key &
  Recovery key access is audited in KMS; dual-control
  required for recovery key retrieval. \\
WAL truncation & Attacker deletes S3 objects &
  S3 versioning enabled; object lock (WORM) for
  compliance-critical buckets. \\
Replay attack & Attacker replays old WAL segments &
  Generation index enforces monotonic WAL positions;
  stale segments rejected. \\
\bottomrule
\end{tabular}
\end{table}

\subsection{SQLCipher Interaction}

When the application database uses SQLCipher (AES-256-CBC at rest),
the replication layer encrypts already-encrypted data with age. This
is deliberate double encryption:

\begin{itemize}[nosep]
  \item SQLCipher protects data at rest on local disk against pod
    filesystem access.
  \item age protects data at rest in S3 against object storage access.
  \item The two keys are independent: SQLCipher key is derived from
    KMS via a different path than the age key.
  \item The replicator does not need the SQLCipher key. It treats the
    database as an opaque byte stream.
\end{itemize}

%% ============================================================
